
This work addressed the question of ranking stability under distribution shift by conducting a systematic Kendall-$\tau$ analysis between ImageNet and ImageNet-V2 leaderboards.

\subsubsection{Summary}

The main contributions are:

\begin{enumerate}
\item \textbf{Ranking stability measurement.} Across 96 models, Kendall-$\tau$ = 0.9647 with 95\% CI [0.9454, 0.9795], substantially above the pre-registered threshold of 0.90.
\end{enumerate}

\begin{enumerate}
\item \textbf{Uniform degradation mechanism.} Accuracy drops are approximately uniform (mean 11.68\%, SD 1.87\%), explaining why rankings remain stable.
\end{enumerate}

\begin{enumerate}
\item \textbf{Practical validation.} Model selection decisions based on ImageNet benchmarks generalize to ImageNet-V2-like distributions.
\end{enumerate}

\subsubsection{Future Directions}

\begin{itemize}
\item \textbf{Testing Alternative Distribution Shifts.} ObjectNet, ImageNet-Sketch, and ImageNet-R may show different ranking stability patterns.
\item \textbf{Architecture-Stratified Analysis.} Different architecture families (e.g., ViT vs. CNN) may exhibit different stability characteristics under distribution shift.
\item \textbf{Theoretical Development.} What properties of a distribution shift determine whether rankings are preserved?
\end{itemize}

\subsubsection{Closing Remarks}

The distinction between accuracy degradation and ranking instability is subtle but practically important. Rankings are preserved ($\tau$ = 0.96) despite accuracy drops (~12\%), which validates the continued use of ImageNet as a model selection benchmark for ImageNet-V2-like distribution shifts.

